% chuong2.tex — sửa nội dung chương ở đây
\section{Vấn đề nghiên cứu}\label{sec:2}

\subsection{Đặt vấn đề}\label{sec:2.1}

\subsubsection{Luật đã cấm nhưng hành vi vẫn diễn ra công khai}\label{sec:2.1.1}

Mua bán thông tin cá nhân của người khác đã bị cấm từ năm 2016: Luật An toàn thông tin mạng 2015 cấm “kinh doanh trái pháp luật thông tin cá nhân của người khác”. Bộ luật Hình sự 2015 cũng coi việc mua bán thông tin riêng của người khác trên mạng mà không được phép là tội phạm (Điều 288). Từ 1/7/2023, Nghị định 13/2023/NĐ-CP nói thẳng hơn: dữ liệu cá nhân không được mua, bán dưới mọi hình thức.

Tuy vậy, taphoammo.net vẫn hoạt động khoảng 4 năm, từ cuối 2021 đến đầu 2026. Sàn không ẩn mình: có tên miền riêng, quảng cáo trên Facebook, TikTok, Google, có công ty đứng sau và hơn 1,35 triệu tài khoản. Khoảng cách giữa “luật đã cấm” và “hành vi vẫn diễn ra” chính là vấn đề mà bài này muốn tìm hiểu.

\subsubsection{Hệ quả đối với người dân}\label{sec:2.1.2}

Hệ quả đầu tiên là quyền riêng tư bị xâm phạm. Người có dữ liệu bị bán thường không biết dữ liệu của mình đã bị bán, bán cho ai, và cũng không có cách nào đòi lại. Một khi đã bị sao chép, dữ liệu có thể được bán lại nhiều lần.

Hệ quả thứ hai là dữ liệu bị bán trở thành nguyên liệu cho các hành vi phạm tội khác. Có thể hình dung một chuỗi đơn giản: kẻ lừa đảo mua một tài khoản Facebook đã có lịch sử hoạt động, mua thêm danh sách số điện thoại kèm họ tên, rồi dùng chúng để mạo danh hoặc nhắn tin dụ dỗ nạn nhân. Tài khoản có sẵn giúp họ vượt qua cơ chế chống tài khoản giả của nền tảng; dữ liệu cá nhân giúp họ tạo lòng tin. Đây là lý do cơ quan điều tra nêu rõ hàng hóa trên taphoammo.net được dùng cho lừa đảo, giả mạo danh tính và tạo tài khoản ảo.

Với hơn 53 triệu giao dịch, số người có thể bị ảnh hưởng là rất lớn, dù hiện chưa có số liệu công bố về số nạn nhân cụ thể.

\subsubsection{Ba câu hỏi nghiên cứu}\label{sec:2.1.3}

Từ vấn đề trên, bài đặt ra ba câu hỏi. Mỗi câu hỏi được trả lời ở một phần cụ thể của bài:

{\small\setstretch{1.1}
\begin{table}[H]
\centering
\caption{Ba câu hỏi nghiên cứu}\label{tab:2.3}
\begin{tabularx}{\textwidth}{>{\raggedright\arraybackslash}l>{\hsize=1.000\hsize\raggedright\arraybackslash}X>{\raggedright\arraybackslash}l}
\toprule
\textbf{\#} & \textbf{Câu hỏi nghiên cứu} & \textbf{Trả lời ở} \\
\midrule
1 & Hoạt động của sàn taphoammo.net vi phạm những quy định nào? & Chương 3 \\
2 & Điều gì khiến việc áp dụng luật chưa ngăn chặn được sớm? & Chương 4 \\
3 & Cần biện pháp gì để phát hiện và ngăn chặn các sàn tương tự? & Chương 5 và 6 \\
\bottomrule
\end{tabularx}
\end{table}}

Câu hỏi 1 xác định “sai ở đâu”. Câu hỏi 2 giải thích “vì sao sai mà vẫn tồn tại lâu”, và đây là câu hỏi trung tâm của bài. Câu hỏi 3 đề xuất cách làm khác đi, dựa trên những thiếu sót tìm được ở câu hỏi 2. Phần kết luận (mục~\ref{sec:6.2}) sẽ quay lại trả lời lần lượt cả ba câu.

\subsection{Cơ sở pháp lý}\label{sec:2.2}

\subsubsection{Nguyên tắc chọn luật để áp dụng}\label{sec:2.2.1}

Một hành vi chỉ bị xét theo quy định có hiệu lực tại thời điểm hành vi xảy ra. Luật ban hành sau không được dùng để xử lý hành vi đã diễn ra trước đó, nếu luật mới bất lợi hơn cho người vi phạm. Vì taphoammo.net hoạt động từ cuối 2021 đến 2/2026, trong khi khung pháp luật thay đổi nhiều lần trong thời gian đó, bài sắp xếp các văn bản theo thời điểm hiệu lực. Cách sắp xếp này cho biết mỗi giai đoạn hoạt động của sàn chịu sự điều chỉnh của những văn bản nào.

Với quy định cấm mua bán dữ liệu cá nhân, có thể chia thành ba giai đoạn:

\begin{enumerate}
\item \textbf{Cuối 2021 đến 30/6/2023.} Chỉ có Luật An toàn thông tin mạng 2015 cấm “kinh doanh trái pháp luật thông tin cá nhân của người khác” (Điều 7 khoản 5), cùng Điều 288 Bộ luật Hình sự. Chưa có văn bản riêng về dữ liệu cá nhân.
\item \textbf{1/7/2023 đến 31/12/2025.} Nghị định 13/2023/NĐ-CP có hiệu lực, lần đầu nêu rõ “dữ liệu cá nhân không được mua, bán dưới mọi hình thức” (Điều 3 khoản 4) và xếp thông tin tài khoản số vào dữ liệu cá nhân cơ bản (Điều 2 khoản 3 điểm k).
\item \textbf{Từ 1/1/2026.} Luật Bảo vệ dữ liệu cá nhân 2025 có hiệu lực, cấm mua bán dữ liệu cá nhân ở cấp luật (Điều 7 khoản 6) và đặt mức phạt tới 10 lần khoản thu từ việc mua bán. Cùng ngày, Nghị định 356/2025/NĐ-CP hướng dẫn luật này thay thế Nghị định 13/2023~\cite{nd356-2025}. Sàn bị khám xét vào tháng 2/2026, nên giai đoạn này chỉ kéo dài khoảng một đến hai tháng hoạt động, và chưa rõ sàn còn hoạt động sau 1/1/2026 hay không.
\end{enumerate}
\begin{figure}[H]
\centering
% Thời gian hiệu lực văn bản so với thời gian sàn hoạt động (1 năm = 2cm, gốc = năm 2021)
\setstretch{1}
\begin{tikzpicture}[x=2cm, y=1cm, font=\footnotesize]
  % trục năm
  \draw[->] (0,0) -- (6.3,0);
  \foreach \n in {2021,...,2027} {
    \draw (\n-2021,0.08) -- (\n-2021,-0.08) node[below] {\n};
  }
  % các thanh
  \newcommand{\hlbar}[5]{%
    \fill[#5] (#1,#3-0.28) rectangle (#2,#3+0.28);
    \node[anchor=east, align=right] at (-0.05,#3) {#4};}
  \hlbar{0.83}{5.08}{4.2}{Sàn hoạt động}{black!70}
  \hlbar{0}{5.5}{3.2}{Luật An toàn\\thông tin mạng 2015}{blue!45}
  \node[anchor=west, font=\scriptsize] at (0.02,3.2) {từ 1/7/2016};
  \hlbar{2.5}{5.0}{2.2}{Nghị định\\13/2023/NĐ-CP}{teal!55}
  \hlbar{5.0}{6.2}{1.2}{Luật Bảo vệ\\dữ liệu cá nhân 2025}{orange!70}
  \draw[->, orange!70, thick] (6.2,1.2) -- (6.35,1.2);
  % khám xét 2/2026
  \draw[dashed, red!70!black, thick] (5.08,0) -- (5.08,4.8) node[above, font=\scriptsize] {Khám xét 2/2026};
\end{tikzpicture}

\caption{Thời gian hiệu lực của các văn bản cấm mua bán dữ liệu cá nhân so với thời gian sàn hoạt động}
\label{fig:2.1}
\end{figure}

Đường nét đứt là thời điểm khám xét (2/2026). Thanh Nghị định 13/2023 và thanh Luật An toàn thông tin mạng 2015 phủ gần hết thời gian sàn hoạt động; Luật Bảo vệ dữ liệu cá nhân 2025 chỉ trùng với khoảng hai tháng cuối.

Phần lớn thời gian hoạt động của sàn (hơn 4 năm) rơi vào hai giai đoạn đầu, khi chưa có mức phạt theo doanh thu.

\subsubsection{Bảng cơ sở pháp lý}\label{sec:2.2.2}

Bảng gồm hai nhóm. \textbf{Nhóm A} là các văn bản trọng tâm về dữ liệu cá nhân và an ninh mạng, dùng cho ba hành vi thuộc nội dung khởi tố và hành vi đặt máy chủ ở nước ngoài (Chương 3). \textbf{Nhóm B} là văn bản bổ trợ về thương mại điện tử, quảng cáo và thanh toán, dùng cho các hành vi về kiểm duyệt người bán, quảng cáo và thanh toán. Trong mỗi nhóm, văn bản xếp theo ngày hiệu lực.

{\small\setstretch{1.1}
\begin{xltabular}{\textwidth}{>{\raggedright\arraybackslash}l>{\hsize=0.773\hsize\raggedright\arraybackslash}X>{\hsize=0.562\hsize\raggedright\arraybackslash}X>{\hsize=2.127\hsize\raggedright\arraybackslash}X>{\hsize=0.538\hsize\raggedright\arraybackslash}X}
\caption{Cơ sở pháp lý áp dụng}\label{tab:2.4}\\
\toprule
\textbf{Nhóm} & \textbf{Văn bản} & \textbf{Điều khoản} & \textbf{Nội dung áp dụng} & \textbf{Hiệu lực} \\
\midrule
\endfirsthead
\toprule
\textbf{Nhóm} & \textbf{Văn bản} & \textbf{Điều khoản} & \textbf{Nội dung áp dụng} & \textbf{Hiệu lực} \\
\midrule
\endhead
\bottomrule
\endlastfoot
A & Luật An toàn thông tin mạng 2015~\cite{luat-attm-2015} (số 86/2015/QH13) & Điều 7 khoản 5 & Cấm thu thập, sử dụng, phát tán, kinh doanh trái pháp luật thông tin cá nhân của người khác & 1/7/2016 đến 30/6/2026 \\
A & Bộ luật Hình sự 2015~\cite{blhs-2015} (sửa đổi 2017) & Điều 288 khoản 1 & Tội mua bán, trao đổi thông tin riêng hợp pháp của người khác trên mạng mà không được phép. Tội danh đã khởi tố & Từ 1/1/2018 \\
A & Luật An ninh mạng 2018~\cite{luat-anm-2018} (số 24/2018/QH14) & Điều 26 khoản 3 & Doanh nghiệp cung cấp dịch vụ trên mạng tại Việt Nam, có thu thập dữ liệu người dùng, phải lưu trữ dữ liệu đó tại Việt Nam & 1/1/2019 đến 30/6/2026 \\
A & Nghị định 53/2022/NĐ-CP~\cite{nd53-2022} & Điều 26, 27 & Hướng dẫn Điều 26 Luật An ninh mạng: doanh nghiệp trong nước phải lưu trữ dữ liệu người dùng tại Việt Nam & Từ 1/10/2022 \\
A & Nghị định 13/2023/NĐ-CP~\cite{nd13-2023} & Điều 3 khoản 4; Điều 2 khoản 3 điểm k & Dữ liệu cá nhân không được mua, bán dưới mọi hình thức; thông tin tài khoản số của cá nhân là dữ liệu cá nhân cơ bản & 1/7/2023 đến 31/12/2025 \\
A & Luật Bảo vệ dữ liệu cá nhân 2025~\cite{luat-bvdlcn-2025} (số 91/2025/QH15) & Điều 7 khoản 6; Điều 8 & Cấm mua bán dữ liệu cá nhân, trừ khi luật quy định khác; phạt hành chính tối đa 10 lần khoản thu từ hành vi mua bán~\cite{luat-bvdlcn-2025} & Từ 1/1/2026 \\
B & Luật Quảng cáo 2012~\cite{luat-qc-2012} (số 16/2012/QH13) & Điều 7 khoản 1; Điều 8 khoản 1 & Hàng hóa, dịch vụ cấm kinh doanh thuộc diện cấm quảng cáo; cấm quảng cáo các sản phẩm đó & Từ 1/1/2013; sửa đổi bởi Luật 75/2025/QH15 \cite{luat-qc-2012} từ 1/1/2026 \\
B & Nghị định 52/2013/NĐ-CP, sửa đổi bởi Nghị định 85/2021/NĐ-CP & Điều 36 (khoản 8, 9, 11), Điều 4 (khoản 1, điểm a) & Trách nhiệm của thương nhân cung cấp dịch vụ sàn giao dịch thương mại điện tử: xử lý hành vi kinh doanh vi phạm trên sàn, hỗ trợ cơ quan nhà nước, liên đới bồi thường~\cite{nd52-2013} & Bản sửa đổi từ 1/1/2022 đến 30/6/2026 \\
B & Nghị định 101/2012/NĐ-CP, sửa đổi bởi Nghị định 80/2016/NĐ-CP & Điều 6 khoản 5, 6 (các hành vi bị cấm trong thanh toán) & Áp dụng cho giai đoạn trước khi Nghị định 52/2024 có hiệu lực & Đến 30/6/2024 \\
B & Nghị định 52/2024/NĐ-CP~\cite{nd52-2024} & Điều 8 khoản 7, 8, 12 & Khoản 7: cấm cung ứng dịch vụ trung gian thanh toán khi chưa được Ngân hàng Nhà nước cấp phép. Khoản 8: cấm dùng dịch vụ thanh toán để kinh doanh trái pháp luật. Khoản 12: cấm hoạt động sai nội dung giấy phép & Từ 1/7/2024 \\
\end{xltabular}}

\subsubsection{Văn bản mới dùng cho phần giải pháp}\label{sec:2.2.3}

Từ 1/7/2026, Luật An ninh mạng 2025 (số 116/2025/QH15, thông qua ngày 10/12/2025) có hiệu lực, hợp nhất và thay thế Luật An ninh mạng 2018 và Luật An toàn thông tin mạng 2015~\cite{luat-anm-2025}. Quy định lưu trữ dữ liệu tại Việt Nam được giữ lại ở Điều 25 khoản 3 của luật mới. Cùng ngày, Luật Thương mại điện tử 2025 (số 122/2025/QH15) có hiệu lực, yêu cầu sàn xác thực người bán trước khi cho mở gian hàng (riêng nghĩa vụ xác thực điện tử danh tính người bán áp dụng từ 1/1/2027~\cite{cuc-tmdt-2026})~\cite{luat-tmdt-2025}.

Cách dùng hai nhóm văn bản trong bài:

\begin{itemize}
\item \textbf{Phân tích vi phạm (Chương 3)} dùng văn bản có hiệu lực tại thời điểm hành vi, tức các văn bản trong bảng trên.
\item \textbf{Đề xuất giải pháp (Chương 5)} dựa trên khung pháp luật hiện hành: Luật Bảo vệ dữ liệu cá nhân 2025 và Nghị định 356/2025, Luật An ninh mạng 2025, Luật Thương mại điện tử 2025.
\end{itemize}
